\chapter{From Steps to the Archimedean Spiral}\label{sec:spirale}

%═══════════════════════════════════════════════════════════════════════════
\section{Geometric Motivation: From the Discrete Step Pattern to the
Continuous Spiral}
%═══════════════════════════════════════════════════════════════════════════

\subsection{Setting and Guiding Question}

In Chapter~\ref{sec:arithmetik} it was shown that the quadratic sequence
\[
  k(n) = \frac{2n^2 + 3n + 1}{6}
\]
assumes integer values precisely for the indices $n \equiv 1, 5 \pmod{6}$.
The resulting discrete point set
\[
  \mathcal{P} = \bigl\{(n,\, k(n)) \mid n \equiv 1, 5 \pmod{6},\ n \geq 1\bigr\}
\]
has a parabolic shape in Cartesian coordinates, since $k(n)$
is a polynomial of degree two in $n$. For the objective of this work --
the embedding of this discrete structure into a geometrically regular,
three-dimensional space -- the Cartesian representation is, however, not
suitable. It admits no natural periodisation, no
angular structure, and it does not make the underlying modular regularity
of the index $n$ geometrically visible.

\begin{tcolorbox}[
  title={Guiding question of this chapter},
  colback=gray!5,
  colframe=gray!50!black,
  enhanced,
  breakable
]
\emph{Does there exist a partition of the difference sequence of the sequence $k(n)$
into blocks of equal length, such that the resulting block sums can be interpreted as circumferences
of concentric circles and the resulting radius sequence
defines an Archimedean spiral with a closed-form pitch constant?
And if so: is this partition uniquely determined?}
\end{tcolorbox}

The answer to both questions, as we shall prove in the following, is
uniquely \emph{yes}.

%─────────────────────────────────────────────────────────────────────────
\subsection{The Difference Sequence of the Integer Values}
\label{subsec:differenzstruktur}

We denote by $(n_k)_{k \geq 1}$ the ascending sequence of all
indices with $n_k \equiv 1, 5 \pmod{6}$:
\[
  n_1 = 1,\quad n_2 = 5,\quad n_3 = 7,\quad n_4 = 11,\quad
  n_5 = 13,\quad n_6 = 17,\quad n_7 = 19,\quad n_8 = 23,\quad \ldots
\]
The corresponding function values (integer values of the sequence) are:
\[
  p_k = k(n_k)\colon\quad 1,\ 11,\ 20,\ 46,\ 63,\ 105,\ 130,\ 188,\ 221,\ldots
\]
We define the \emph{first difference sequence} of this value sequence:
\[
  \delta_k = p_{k+1} - p_k = k(n_{k+1}) - k(n_k),\qquad k \geq 1.
\]
The first terms of this difference sequence are:
\[
  \delta_1 = 10,\quad \delta_2 = 9,\quad \delta_3 = 26,\quad
  \delta_4 = 17,\quad \delta_5 = 42,\quad \delta_6 = 25,\quad
  \delta_7 = 58,\quad \delta_8 = 33,\quad \ldots
\]

\begin{bemerkung}[Alternating growth of the differences]
The difference sequence $(\delta_k)$ displays a clear alternating pattern:
\begin{itemize}
  \item The terms at \emph{odd} positions ($\delta_1, \delta_3, \delta_5, \ldots$)
        take the values $10, 26, 42, 58, \ldots$ and form an arithmetic subsequence
        with constant second difference $16$.
  \item The terms at \emph{even} positions ($\delta_2, \delta_4, \delta_6, \ldots$)
        take the values $9, 17, 25, 33, \ldots$ and form an arithmetic subsequence
        with constant second difference $8$.
\end{itemize}
This pattern directly reflects the alternation of the index differences $n_{k+1} - n_k$,
which switch between $\Delta n = 4$ (for $n_k \equiv 1 \pmod{6}$, \emph{major steps})
and $\Delta n = 2$ (for $n_k \equiv 5 \pmod{6}$, \emph{minor steps}).
\end{bemerkung}

%─────────────────────────────────────────────────────────────────────────
\subsection{The Grouping Problem and Its Unique Solution}
\label{subsec:gruppenbildung}

We pose the following algebraic question: which partition of the
difference sequence $(\delta_k)$ into consecutive blocks of equal length
$\ell \in \{1, 2, 3, \ldots\}$ generates a sequence of block sums
that forms an \emph{arithmetic progression}?

Formally, for a block size $\ell \geq 1$ we define the sequence of
block sums
\[
  S_m^{(\ell)} = \sum_{j=1}^{\ell} \delta_{\ell(m-1)+j},\qquad m \geq 1,
\]
and ask: for which values of $\ell$ is $(S_m^{(\ell)})_{m \geq 1}$
an arithmetic sequence?

\begin{satz}[Uniqueness of the arithmetic grouping]\label{satz:eindeutigkeit}
Among all block sizes $\ell \in \{1, 2, 3\}$, $\ell = 2$ is the unique value
for which $(S_m^{(\ell)})$ forms an arithmetic progression with non-trivial
common difference.
\end{satz}

\begin{proof}
We examine the three relevant cases.

\textbf{Case $\ell = 1$:} The block sums are the differences themselves:
\[
  S_1^{(1)} = 10,\quad S_2^{(1)} = 9,\quad S_3^{(1)} = 26,\quad
  S_4^{(1)} = 17,\quad \ldots
\]
The sequence is not monotone ($10, 9, 26, 17, \ldots$) and hence not
an arithmetic progression.

\textbf{Case $\ell = 2$:} The block sums are pairwise sums:
\begin{align*}
S_1^{(2)} &= \delta_1 + \delta_2 = 10 + 9 = 19,\\
S_2^{(2)} &= \delta_3 + \delta_4 = 26 + 17 = 43,\\
S_3^{(2)} &= \delta_5 + \delta_6 = 42 + 25 = 67,\\
S_4^{(2)} &= \delta_7 + \delta_8 = 58 + 33 = 91.
\end{align*}
The differences of consecutive block sums are:
\[
  43 - 19 = 24,\qquad 67 - 43 = 24,\qquad 91 - 67 = 24.
\]
For the general proof we use the closed-form expressions derived in
Section~\ref{sec:minor_major_herleitung}
(Theorem~\ref{satz:minor_major_formeln}):
\[
  \delta_{2m-1} = \Delta_{\mathrm{Major}}(m) = 16m - 6,\qquad
  \delta_{2m} = \Delta_{\mathrm{Minor}}(m) = 8m + 1.
\]
This yields:
\[
  S_m^{(2)} = \delta_{2m-1} + \delta_{2m}
            = (16m-6) + (8m+1)
            = 24m - 5,
\]
and
\[
  S_{m+1}^{(2)} - S_m^{(2)}
  = \bigl[24(m+1)-5\bigr] - \bigl[24m-5\bigr]
  = 24.
\]
The sequence $(S_m^{(2)})$ is therefore an arithmetic progression with common
difference $24$.

\textbf{Case $\ell = 3$:} The block sums are:
\begin{align*}
S_1^{(3)} &= \delta_1 + \delta_2 + \delta_3 = 10 + 9 + 26 = 45,\\
S_2^{(3)} &= \delta_4 + \delta_5 + \delta_6 = 17 + 42 + 25 = 84,\\
S_3^{(3)} &= \delta_7 + \delta_8 + \delta_9 = 58 + 33 + 74 = 165.
\end{align*}
The differences: $84 - 45 = 39$ and $165 - 84 = 81$ are distinct.
Hence $(S_m^{(3)})$ is not an arithmetic progression.
\end{proof}

\begin{bemerkung}[Connection to OEIS]
	The sequence of circle circumferences $U_m = 24m - 5$ is listed in the OEIS
	under A350051~\cite{OEIS} as part of the trisection
	of $A017101 = \{3 + 8k\}_{k \geq 0}$. It is congruent to
	$1 \pmod{6}$ -- in agreement with the congruence structure
	of the sequence $k(n)$.
\end{bemerkung}

\begin{anmerkung}[Extension to all $\ell \geq 1$]
Theorem~\ref{satz:eindeutigkeit} explicitly treats $\ell \in \{1,2,3\}$.
For general $\ell \geq 1$ one can directly verify that
\[
  S_m^{(\ell)} = \sum_{j=1}^{\ell} \delta_{\ell(m-1)+j}
\]
forms an arithmetic progression in $m$ if and only if
$S_{m+1}^{(\ell)} - S_m^{(\ell)}$ is independent of $m$.
Since the difference sequence $(\delta_j)$ has period~$2$
(in the sense of the alternating minor/major structure),
the difference $S_{m+1}^{(\ell)} - S_m^{(\ell)}$ is
$m$-independent only when $\ell$ is a multiple of~$2$.
For odd $\ell$, each block contains a different mixture
of minor and major steps, leading to $m$-dependent differences.
Among all even $\ell$, $\ell = 2$ yields the minimal and
geometrically natural pairing. A complete proof for all
$\ell \geq 4$ via a general parity argument remains as
a supplement.
\end{anmerkung}

%─────────────────────────────────────────────────────────────────────────
\subsection{Geometric Interpretation of the Block Sums as Circle Circumferences}
\label{subsec:geometrische_interpretation}

Having established that the pairwise block sums $S_m^{(2)} = 24m - 5$ form an arithmetic
progression, we now present the decisive geometric argument.

An Archimedean spiral $r(\varphi) = a\varphi$ has the defining
property of \emph{constant radius difference}: for each complete
revolution $\Delta\varphi = 2\pi$:
\[
  \Delta r = 2\pi a = \text{const}.
\]
Equivalently: the circumferences $U_m = 2\pi R_m$ of consecutive
spiral windings themselves form an arithmetic progression with constant
difference $\Delta U = 2\pi \Delta R$.
\clearpage

\begin{tcolorbox}[
  title={Justification of the circumference identification},
  colback=blue!3,
  colframe=blue!40!black,
  enhanced,
  breakable
]
The integer points of the sequence $k(n)$ exhibit, in a polar representation,
a radial arrangement: points with $n_k \equiv 1 \pmod{6}$ cluster
along a preferred radial direction, points with $n_k \equiv 5 \pmod{6}$
along another. This structure suggests treating each pair of two consecutive
points -- one major step and one minor step -- as a complete
\emph{period} of the arrangement, geometrically corresponding to one winding
of the spiral.

Under this interpretation, the sum of the radius increments
within such a period corresponds to the circumference increment of a circle. The
fact that this sum -- and only this pairing -- yields an arithmetic
progression (Theorem~\ref{satz:eindeutigkeit}) constitutes the algebraic
proof that the identification is geometrically consistent: only with this
pairing does a spiral with \emph{constant} winding separation emerge.
\end{tcolorbox}

\begin{definition}[Circle circumferences of the sequence]\label{def:kreisumfaenge}
We identify the pairwise block sums of the difference sequence
with the circumferences of a family of concentric circles:
\[
  U_m \coloneqq S_m^{(2)} = \delta_{2m-1} + \delta_{2m} = 24m - 5.
\]
\end{definition}

%═══════════════════════════════════════════════════════════════════════════
\section{Derivation of the Major and Minor Step Formulae}
\label{sec:minor_major_herleitung}
%═══════════════════════════════════════════════════════════════════════════

\subsection{Explicit Computation of the Step Sizes}

We derive the closed-form expressions for the major and minor step sizes
directly from the definition of $k(n)$.

The indices of the integer values follow a periodic pattern:
\[
  \ldots,\quad 6m-5,\quad 6m-1,\quad 6m+1,\quad 6m+5,\quad 6m+7,\quad\ldots
\]
Here $6m-5 \equiv 1 \pmod{6}$ and $6m-1 \equiv 5 \pmod{6}$.
The period $m$ ($m = 1, 2, 3, \ldots$) consists of:
\begin{itemize}
  \item a \textbf{major step} from $n = 6m-5$ to $n = 6m-1$
        (index difference $\Delta n = 4$),
  \item a \textbf{minor step} from $n = 6m-1$ to $n = 6m+1$
        (index difference $\Delta n = 2$).
\end{itemize}

\begin{satz}[Explicit formulae for major and minor steps]\label{satz:minor_major_formeln}
For the $m$-th period ($m \geq 1$):
\begin{align}
  \Delta_{\mathrm{Major}}(m)
  &= k(6m-1) - k(6m-5)
  = 16m - 6,\label{eq:major_formel}\\[6pt]
  \Delta_{\mathrm{Minor}}(m)
  &= k(6m+1) - k(6m-1)
  = 8m + 1.\label{eq:minor_formel}
\end{align}
\end{satz}

\begin{proof}
\textbf{Major step:} We compute $k(6m-1) - k(6m-5)$ directly.
With $k(n) = \frac{2n^2+3n+1}{6}$:
\begin{align*}
  k(6m-1) - k(6m-5)
  &= \frac{2(6m-1)^2 + 3(6m-1) + 1}{6}
   - \frac{2(6m-5)^2 + 3(6m-5) + 1}{6} \\[4pt]
  &= \frac{1}{6}\bigl[2\bigl((6m-1)^2 - (6m-5)^2\bigr)
     + 3\bigl((6m-1) - (6m-5)\bigr)\bigr].
\end{align*}
We use the difference of squares identity $a^2 - b^2 = (a+b)(a-b)$
with $a = 6m-1$ and $b = 6m-5$:
\begin{align*}
  (6m-1)^2 - (6m-5)^2
  &= \bigl[(6m-1)+(6m-5)\bigr]\bigl[(6m-1)-(6m-5)\bigr] \\
  &= (12m-6)\cdot 4
  = 48m - 24.
\end{align*}
Furthermore: $(6m-1)-(6m-5) = 4$. Substituting:
\begin{align*}
  k(6m-1) - k(6m-5)
  &= \frac{1}{6}\bigl[2(48m-24) + 3 \cdot 4\bigr] \\
  &= \frac{1}{6}\bigl[96m - 48 + 12\bigr] \\
  &= \frac{96m - 36}{6}
  = 16m - 6. \checkmark
\end{align*}

\textbf{Minor step:} We compute $k(6m+1) - k(6m-1)$:
\begin{align*}
  k(6m+1) - k(6m-1)
  &= \frac{1}{6}\bigl[2\bigl((6m+1)^2 - (6m-1)^2\bigr)
     + 3\bigl((6m+1) - (6m-1)\bigr)\bigr].
\end{align*}
With $a = 6m+1$, $b = 6m-1$:
\begin{align*}
  (6m+1)^2 - (6m-1)^2
  &= (12m)(2) = 24m.
\end{align*}
Furthermore: $(6m+1)-(6m-1) = 2$. Substituting:
\begin{align*}
  k(6m+1) - k(6m-1)
  &= \frac{1}{6}\bigl[2 \cdot 24m + 3 \cdot 2\bigr]
  = \frac{48m + 6}{6}
  = 8m + 1. \checkmark
\end{align*}
\end{proof}

\begin{table}[htbp]
\centering
\caption{Step sizes for the first periods}
\begin{tabular}{c|cc|cc|c}
\toprule
Period $m$ & Starting value $n$ & $\Delta_{\mathrm{Major}}$ &
Intermediate value $n$ & $\Delta_{\mathrm{Minor}}$ & $U_m$ \\
\midrule
1 & 1  & 10 & 5  & 9  & 19 \\
2 & 7  & 26 & 11 & 17 & 43 \\
3 & 13 & 42 & 17 & 25 & 67 \\
4 & 19 & 58 & 23 & 33 & 91 \\
5 & 25 & 74 & 29 & 41 & 115 \\
\bottomrule
\end{tabular}
\end{table}

\begin{bemerkung}[Verification of the formulae]
For $m=1$: $\Delta_{\mathrm{Major}}(1) = 16\cdot1 - 6 = 10$ ,
$\Delta_{\mathrm{Minor}}(1) = 8\cdot1 + 1 = 9$ .
For $m=2$: $\Delta_{\mathrm{Major}}(2) = 32-6 = 26$ ,
$\Delta_{\mathrm{Minor}}(2) = 17$ . The formulae agree
exactly with the computed differences.
\end{bemerkung}

\subsection{The Unified Generator Function}

The two subsequences can be combined into a single formula
that generates all steps of the interleaved sequence $(\delta_k)$.

\begin{satz}[Generating function of the difference sequence]
The complete difference sequence $(\delta_k)_{k \geq 1}$ in its natural
order starting from $n = 1$:
\[
  \delta_1 = 10,\; \delta_2 = 9,\; \delta_3 = 26,\; \delta_4 = 17,\;
  \delta_5 = 42,\; \delta_6 = 25,\; \delta_7 = 58,\; \delta_8 = 33,\; \ldots
\]
(\textbf{major, minor, major, minor, \ldots}) is generated by the closed-form expression
\begin{equation}
  a_k = \frac{1}{2}\bigl[(12k+3) + (-1)^{k+1}(4k+1)\bigr],
  \qquad k \geq 1,
\end{equation}
where $a_1 = 10$ corresponds to the first major step (from $n_1 = 1$ to $n_2 = 5$).
\end{satz}

\begin{proof}
For $k = 2m-1$ (odd positions, major steps):
\begin{align*}
  a_{2m-1}
  &= \frac{1}{2}\bigl[(12(2m-1)+3) + (-1)^{2m}(4(2m-1)+1)\bigr] \\
  &= \frac{1}{2}\bigl[(24m-9) + (8m-3)\bigr]
  = \frac{32m-12}{2} = 16m-6.
\end{align*}
This is precisely $\Delta_{\mathrm{Major}}(m) = 16m-6$. \quad
Verification: $a_1 = 16\cdot1-6 = 10$. \checkmark

For $k = 2m$ (even positions, minor steps):
\begin{align*}
  a_{2m}
  &= \frac{1}{2}\bigl[(12(2m)+3) + (-1)^{2m+1}(4(2m)+1)\bigr] \\
  &= \frac{1}{2}\bigl[(24m+3) - (8m+1)\bigr]
  = \frac{16m+2}{2} = 8m+1.
\end{align*}
This is precisely $\Delta_{\mathrm{Minor}}(m) = 8m+1$. \quad
Verification: $a_2 = 8\cdot1+1 = 9$. \checkmark

The period sums remain identical:
\[
  U_m = a_{2m-1} + a_{2m}
      = (16m-6) + (8m+1)
      = 24m - 5. \qquad\checkmark
\]
\end{proof}

%═══════════════════════════════════════════════════════════════════════════
\section{From Steps to Circle Circumferences}
%═══════════════════════════════════════════════════════════════════════════

With the formulae proved in Section~\ref{sec:minor_major_herleitung},
we can now introduce the definitions of the minor and major steps.

\begin{definition}[Minor and Major Steps]
The step sizes of the $m$-th period are denoted as follows:
\begin{align}
  \Delta_{\mathrm{Minor},m} &= 8m + 1 \quad
    \text{(step $\Delta n = 2$, from } n\equiv5 \text{ to } n\equiv1\text{)},\\
  \Delta_{\mathrm{Major},m} &= 16m - 6 \quad
    \text{(step $\Delta n = 4$, from } n\equiv1 \text{ to } n\equiv5\text{)}.
\end{align}
\end{definition}

\begin{definition}[Circle circumferences]
One major and one minor step together form the $m$-th
\emph{circle circumference}:
\begin{equation}
  U_m = \Delta_{\mathrm{Major},m} + \Delta_{\mathrm{Minor},m}.
\end{equation}
\end{definition}

\begin{satz}[Arithmetic sequence of circumferences]
The circle circumferences form an arithmetic sequence:
\begin{equation}
  U_m = \Delta_{\mathrm{Major},m} + \Delta_{\mathrm{Minor},m}
      = (16m-6) + (8m+1)
      = 24m - 5,
\end{equation}
with constant difference
\[
  \Delta U = U_{m+1} - U_m = 24.
\]
\end{satz}

%═══════════════════════════════════════════════════════════════════════════
\section{Radius Sequence}
%═══════════════════════════════════════════════════════════════════════════

From the circumferences, the radius sequence follows via $U_m = 2\pi R_m$:
\begin{equation}
  R_m = \frac{U_m}{2\pi} = \frac{24m - 5}{2\pi}.
\end{equation}

\begin{satz}[Constant radius difference -- key result]
\begin{equation}
  \boxed{\Delta R = R_{m+1} - R_m = \frac{24}{2\pi} = \frac{12}{\pi}
  \approx 3.819718634205}
\end{equation}
\end{satz}

\begin{proof}
Direct computation:
\[
  R_{m+1} - R_m
  = \frac{24(m+1)-5}{2\pi} - \frac{24m-5}{2\pi}
  = \frac{24}{2\pi}
  = \frac{12}{\pi}.
\]
\end{proof}

%═══════════════════════════════════════════════════════════════════════════
\section{Archimedean Spiral}
%═══════════════════════════════════════════════════════════════════════════

\begin{definition}
An \emph{Archimedean spiral} is a plane curve in polar coordinates
of the form
\[
  r(\theta) = a\theta,
\]
with pitch constant $a > 0$. The characterising property of this
class of curves is the constant radius difference between consecutive
windings: for each complete revolution $\Delta\theta = 2\pi$:
\[
  \Delta r = r(\theta + 2\pi) - r(\theta) = 2\pi a.
\]
\end{definition}

The uniqueness of the pairing proved in Theorem~\ref{satz:eindeutigkeit}
and the resulting constant $\Delta R = 12/\pi$ determine the
pitch constant $a$ uniquely:
\begin{align}
  2\pi a &= \Delta R = \frac{12}{\pi},\\[4pt]
  a &= \frac{12}{2\pi^2} = \frac{6}{\pi^2}.
\end{align}

\begin{satz}[Spiral equation]
\begin{equation}
  \boxed{
    r(\theta) = \frac{6}{\pi^2}\,\theta
    \approx 0.607927101854 \cdot \theta
  }
\end{equation}
\end{satz}

\begin{bemerkung}[Connection to the Basel problem]
The pitch constant
\[
  a = \frac{6}{\pi^2} = \frac{1}{\zeta(2)},
  \qquad
  \zeta(2) = \sum_{n=1}^{\infty} \frac{1}{n^2} = \frac{\pi^2}{6}
\]
is the reciprocal of the Riemann zeta function at $s=2$, i.e.\
the solution to the classical Basel problem. The value $\zeta(2) = \pi^2/6$
was proved by Leonhard Euler in 1734~\cite{Euler1740,Hardy2008}. The fact that the spiral constant,
obtained from purely combinatorial-arithmetic considerations on the
difference structure of the sequence $k(n)$, coincides with $1/\zeta(2)$
is a remarkable structural connection.
\end{bemerkung}
